\section{Designing Interactive Systems: Methods}

\begin{ejercicio}
    You have been tasked with designing a social robot for a healthcare facility. How would you identify and involve the relevant stakeholders in the design process? Name two methods that help to involve stakeholders and explain also how these methods could be applied.

    In order to identify and involve relevant stakeholders in the design process of a social robot for a healthcare facility, I would first conduct a stakeholder analysis to identify key groups such as healthcare professionals (doctors, nurses), patients, caregivers, and administrative staff. Specifically, I would use ``Personas'' to create detailed profiles of representative users. Regarding methods to involve stakeholders, I would use:
    \begin{enumerate}
        \item \ul{Interviews}: Conducting one-on-one interviews with stakeholders to gather their insights, needs, and concerns regarding the social robot. This method allows for in-depth understanding of their perspectives and can help identify specific requirements for the robot's design.
        \item \ul{Surveys}: Distributing surveys to a larger group of stakeholders to collect quantitative data on their preferences, expectations, and potential concerns. This method can help identify common trends and priorities among stakeholders, which can inform the design process.
    \end{enumerate}
\end{ejercicio}

\begin{ejercicio}
    You  are researching the use of AI in the classroom. A chatbot with generative AI is to be used in the classroom by teachers as well as students. In the case of the stakeholder analysis, create at least two personas. Which personas that would be? Which kind of stakeholders?
    \begin{enumerate}
        \item \ul{Persona 1: Teacher Persona}
        
        Sarah is a 55-year-old teacher who does not usually use technology in her classroom. She is quite stressed due to the big number of students she has to manage and the limited time she has to prepare her lessons. She is interested in using the chatbot to help her with lesson planning and to provide additional support to her students, but she is concerned about the reliability and accuracy of the AI-generated content.

        \item \ul{Persona 2: Student Persona}
        
        Alex is a 16-year-old high school student who is tech-savvy and enjoys using digital tools for learning. He is interested in using the chatbot to get instant answers to his questions and to receive personalized learning recommendations. However, he is concerned about the privacy of his data and the potential for the AI to provide biased or incorrect information.
    \end{enumerate}
\end{ejercicio}

\begin{ejercicio}
    Name three characteristics of prototype and explain two differences between a prototype and a mock-up.

    Prototypes are early models of a product that are used to test and validate design concepts. Three characteristics of prototypes are:
    \begin{enumerate}
        \item \ul{Functionality}: Prototypes often have some level of functionality that allows users to interact with the product and test its features.
        \item \ul{Iterative}: Prototypes are typically developed in an iterative process, allowing for continuous refinement and improvement based on user feedback.
        \item \ul{Representation}: Prototypes can represent different aspects of the product, such as its appearance, behavior, or user interface, depending on the stage of development.
    \end{enumerate}

    Two differences between a prototype and a mock-up are:
    \begin{enumerate}
        \item \ul{Interactivity}: Prototypes are interactive and allow users to test the functionality of the product, while mock-ups are static representations that focus on the visual design and layout without providing any interactive features.
        \item \ul{Purpose}: Prototypes are used to validate design concepts and gather user feedback, while mock-ups are primarily used to communicate design ideas and aesthetics to stakeholders.
    \end{enumerate}
\end{ejercicio}

\begin{ejercicio}
    Briefly describe the Wizard of Oz method and why it is needed. Give an example of a scenario where this method could be applied.

    The Wizard of Oz method is a type of prototyping technique where users interact with a system that they believe to be autonomous, but is actually being operated or controlled by a human behind the scenes. This method is needed to test and evaluate user interactions with a system that may not yet be fully developed or functional, allowing designers to gather feedback and insights without the need for a complete implementation.
\end{ejercicio}

\begin{ejercicio}
    What is the Double Diamond model? Explain the two diamonds and the phases of the design process according to this model.

    The Double Diamond model is a design thinking framework that consists of two diamonds, each representing a different phase of the design process. The first diamond represents the definition of the problem, and consists of two phases: \textit{Discover} and \textit{Define}. The second diamond represents the development of the solution, and consists of two phases: \textit{Develop} and \textit{Deliver}. The model emphasizes the importance of divergent and convergent thinking in both phases, allowing designers to explore a wide range of ideas before narrowing down to a final solution.
\end{ejercicio}